\chapter{Proposed Methodology}
\label{ch:methodology}

\section{Overview}
\label{sec:method_overview}

This chapter details the proposed experimental methodology designed to rigorously evaluate and quantify the impact of experimental randomness on machine learning models. The framework isolates random seed variations while preserving consistent feature representations, algorithmic hyperparameters, and evaluation protocols. By comparing Decision Tree, Random Forest, and Linear Support Vector Machine on the Forest CoverType dataset, the methodology evaluates both point classification efficacy and statistical distribution stability.

\section{System Requirements}
\label{sec:sys_req}

To ensure complete computational reproducibility, the hardware and software environments were strictly standardized as detailed in Table \ref{tab:system_requirements}.

\begin{table}[H]
\centering
\caption{Computational Environment and System Specifications}
\label{tab:system_requirements}
\begin{tabular}{p{4cm} p{10cm}}
\toprule
\textbf{Component} & \textbf{Specification / Value} \\
\midrule
Operating Environment & Google Colaboratory / Linux x86\_64 (Ubuntu 22.04 LTS) \\
Host Architecture & Intel Xeon CPU @ 2.20GHz (2 vCPUs) \\
Physical Memory (RAM) & 12.67 GB Available System RAM \\
Disk Storage & 107.72 GB Available Virtual Disk \\
Programming Language & Python 3.10.12 \\
Core Libraries & Scikit-learn (1.2.2+), NumPy (1.24.3+), Pandas (2.0.3+), Matplotlib (3.7.1+), Seaborn (0.12.2+) \\
Development Interface & Jupyter Notebook (.ipynb) \\
\bottomrule
\end{tabular}
\end{table}

\section{Proposed Approach}
\label{sec:proposed_approach}

The experimental architecture is structured into sequential investigatory phases:
\begin{enumerate}
    \item \textbf{Data Ingestion \& Exploratory Analysis:} Load the 581,012-sample dataset, audit class distributions, verify absence of missing values, and construct an exact 10,000-sample stratified subset.
    \item \textbf{Phase 1 --- Baseline Modeling:} Split data into 80\% training and 20\% testing sets using a fixed baseline seed ($S_0 = 42$), fit each classifier, and compute multi-class evaluation metrics.
    \item \textbf{Phase 2 --- Split Seed Sensitivity Analysis:} Repeat training and evaluation across five distinct random split seeds:
    \begin{equation}
        \mathcal{S}_{\text{split}} = \{0, 1, 42, 99, 123\}
    \end{equation}
    Measure empirical mean $\mu$ and standard deviation $\sigma$ across runs to capture estimator variance.
    \item \textbf{Phase 3 --- Cross-Validation \& Reproducibility Verification:} Perform 5-fold stratified cross-validation across three shuffle seeds ($\mathcal{S}_{\text{cv}} = \{11, 22, 33\}$) and execute duplicate back-to-back runs to verify bitwise deterministic output.
    \item \textbf{Phase 4 --- Dataset Size Scaling Experiment:} Evaluate model learning dynamics across three training sample volumes ($N \in \{2000, 5000, 8000\}$) on a standardized holdout test set to quantify sample complexity.
    \item \textbf{Phase 5 --- Hyperparameter Optimization Analysis:} Benchmark inductive bias variations by tuning tree depth ($d \in \{5, 10, \text{None}\}$) for Decision Trees and ensemble size ($B \in \{50, 100, 200\}$) for Random Forests.
    \item \textbf{Phase 6 --- Feature Group Contribution Analysis:} Perform feature ablation comparing the 10 continuous terrain features against the complete 54-feature set to assess dimensional utility.
\end{enumerate}

\section{Algorithms and Computational Models}
\label{sec:algorithms}

Three diverse machine learning paradigms were selected to evaluate how model architecture influences sensitivity to stochastic perturbations.

\subsection{Decision Tree Classifier (CART)}
Classification and Regression Tree (CART) recursively partitions the input space $\mathcal{X}$ into axis-aligned rectangular regions. At each internal node $m$, the algorithm identifies the optimal feature $j$ and split threshold $t$ that minimizes node impurity. Using the Gini Impurity index:
\begin{equation}
    I_G(m) = 1 - \sum_{k=1}^{K} p_{mk}^2
\end{equation}
where $p_{mk}$ represents the proportion of training instances belonging to class $k$ in node $m$. The split selection solves:
\begin{equation}
    \arg\min_{j, t} \left[ \frac{N_L}{N_m} I_G(m_L) + \frac{N_R}{N_m} I_G(m_R) \right]
\end{equation}
Decision trees are prone to high variance: minor shifts in sample distributions alter early split decisions, cascading through subtrees.

\subsection{Random Forest Classifier}
Random Forest constructs an ensemble of $B$ de-correlated decision trees $\{T_b\}_{b=1}^B$. For each tree $b$:
\begin{enumerate}
    \item A bootstrap sample $\mathcal{D}_b$ of size $N$ is drawn uniformly with replacement from the training set $\mathcal{D}$;
    \item At each node, a random subset of $m \approx \sqrt{p}$ features is selected from the total $p$ features to determine the split;
    \item Ensemble prediction is obtained via majority voting across trees:
    \begin{equation}
        \hat{y} = \arg\max_k \sum_{b=1}^{B} \mathbb{I}(T_b(\mathbf{x}) = k)
    \end{equation}
\end{enumerate}
By averaging across $B$ de-correlated estimators, ensemble variance is reduced by a factor proportional to the correlation between individual trees:
\begin{equation}
    \text{Var}(\bar{T}) = \rho \sigma^2 + \frac{1 - \rho}{B} \sigma^2
\end{equation}
where $\rho$ is tree-pair correlation and $\sigma^2$ is individual tree variance.

\subsection{Linear Support Vector Machine (Linear SVM)}
Linear Support Vector Classification optimizes a maximum-margin hyperplane separating classes in feature space. For multiclass classification with $K$ classes, a One-vs-Rest (OvR) formulation is utilized. For each binary subproblem, the objective minimizes the regularized soft-margin hinge loss:
\begin{equation}
    \min_{\mathbf{w}, b, \boldsymbol{\xi}} \frac{1}{2} \|\mathbf{w}\|^2 + C \sum_{i=1}^N \xi_i \quad \text{s.t.} \quad y_i (\mathbf{w}^T \mathbf{x}_i + b) \ge 1 - \xi_i, \quad \xi_i \ge 0
\end{equation}
where $C > 0$ governs the trade-off between margin maximization and margin violation penalties, and $\xi_i$ are slack variables. Because optimization depends on inner products, standardizing continuous features is critical for convergence.

\section{Data Collection and Preprocessing}
\label{sec:data_preprocessing}

\subsection{Dataset Description}
The Forest CoverType dataset comprises 581,012 geological survey records from four wilderness areas in Roosevelt National Forest, Colorado. The objective is to predict the principal forest cover type (Classes 1 through 7):
\begin{itemize}
    \item Class 1: Spruce/Fir
    \item Class 2: Lodgepole Pine
    \item Class 3: Ponderosa Pine
    \item Class 4: Cottonwood/Willow
    \item Class 5: Aspen
    \item Class 6: Douglas-fir
    \item Class 7: Krummholz
\end{itemize}
The 54 input features encompass:
\begin{itemize}
    \item 10 Continuous features: Elevation, Aspect, Slope, Horizontal/Vertical Distance to Hydrology, Horizontal Distance to Roadways, Hillshade (9am, Noon, 3pm), Horizontal Distance to Fire Points.
    \item 4 Binary Wilderness Area indicators (Wilderness Area 1 to 4).
    \item 40 Binary Soil Type indicators (Soil Type 1 to 40).
\end{itemize}

\subsection{Stratified Subsampling}
Training iterative multi-seed experiments on the full 581,012 samples requires extensive computational time. To ensure rigorous multi-run evaluation while maintaining computational feasibility, a stratified random sample of $N = 10,000$ instances was generated using \texttt{StratifiedShuffleSplit} with a fixed seed ($S_{\text{sample}} = 42$). Table \ref{tab:class_distribution} demonstrates that the sample accurately mirrors the parent population distribution.

\begin{table}[H]
\centering
\caption{Forest CoverType Class Distribution Comparison}
\label{tab:class_distribution}
\begin{tabular}{clrrrr}
\toprule
\textbf{Class} & \textbf{Cover Type Name} & \textbf{Full Count} & \textbf{Full \%} & \textbf{Sample Count} & \textbf{Sample \%} \\
\midrule
1 & Spruce/Fir & 211,840 & 36.46\% & 3,646 & 36.46\% \\
2 & Lodgepole Pine & 283,301 & 48.76\% & 4,876 & 48.76\% \\
3 & Ponderosa Pine & 35,754 & 6.15\% & 615 & 6.15\% \\
4 & Cottonwood/Willow & 2,747 & 0.47\% & 47 & 0.47\% \\
5 & Aspen & 9,493 & 1.63\% & 163 & 1.63\% \\
6 & Douglas-fir & 17,367 & 2.99\% & 299 & 2.99\% \\
7 & Krummholz & 20,510 & 3.53\% & 353 & 3.53\% \\
\midrule
\textbf{Total} & & \textbf{581,012} & \textbf{100.00\%} & \textbf{10,000} & \textbf{100.00\%} \\
\bottomrule
\end{tabular}
\end{table}

\subsection{Data Hygiene and Feature Normalization}
Data preprocessing adhered to rigorous standards to prevent data leakage:
\begin{enumerate}
    \item \textbf{Leakage Prevention:} Feature scalers were fitted strictly on the training partition ($\mathbf{X}_{\text{train}}$) and subsequently applied to transform test instances ($\mathbf{X}_{\text{test}}$).
    \item \textbf{Continuous Feature Standardization:} The 10 continuous cartographic variables were normalized using z-score standardization:
    \begin{equation}
        z = \frac{x - \mu_{\text{train}}}{\sigma_{\text{train}}}
    \end{equation}
    \item \textbf{Preservation of Binary Encodings:} The 44 one-hot encoded binary variables (Wilderness Areas and Soil Types) were preserved in raw boolean format $\{0, 1\}$ to maintain sparsity and prevent artificial distortion of categorical distances.
\end{enumerate}

\section{Theoretical Formulations of Extended Experiments}
\label{sec:extended_exp_theory}

To provide an exhaustive characterization of model behavior beyond basic random seed stability, three targeted experimental extensions were formulated:

\subsection{Theoretical Foundations of Dataset Size Scaling}
\label{sec:theory_data_size}

The volume of available training data directly influences learning capacity and generalization error. In statistical learning theory, the generalization error bound of an estimator $h \in \mathcal{H}$ with Vapnik-Chervonenkis (VC) dimension $d_{\text{VC}}$ over $N$ independent training samples scales according to:
\begin{equation}
    \epsilon(h) \le \hat{\epsilon}(h) + \mathcal{O}\left( \sqrt{\frac{d_{\text{VC}} \log(N / d_{\text{VC}}) + \log(1/\delta)}{N}} \right)
\end{equation}
with probability at least $1 - \delta$. As sample volume $N$ increases, the variance term decreases asymptotically, shrinking the gap between empirical training error and true generalization risk. 

In multi-class imbalanced ecological datasets such as Forest CoverType, dataset size plays a critical role in minority class boundary estimation. With rare classes comprising as little as $0.47\%$ of the population (Class 4, Cottonwood/Willow), a small training sample of $N=2,000$ provides fewer than 10 training instances of rare classes, starving split criteria of statistical power. Testing across $N \in \{2000, 5000, 8000\}$ systematically evaluates sample complexity, empirical learning rate, and minority-class sensitivity on a held-out test set ($N_{\text{test}} = 2,000$).

\subsection{Theoretical Principles of Hyperparameter Tuning}
\label{sec:theory_hyperparameter}

Hyperparameters are external algorithmic configurations chosen prior to training that govern model flexibility, hypothesis space complexity, and inductive bias. Unlike internal parameters learned via optimization (such as tree split thresholds or SVM weights), hyperparameters control the regularization-capacity trade-off:
\begin{itemize}
    \item \textbf{Decision Tree Maximum Depth (\texttt{max\_depth}):} Controls the maximum allowable length of root-to-leaf paths. An unconstrained tree (\texttt{max\_depth=None}) recursively partitions data until leaves are pure ($I_G = 0$). While this achieves zero training error, it models noise and localized anomalies, leading to severe overfitting. Restricting depth to $d \in \{5, 10\}$ acts as structural regularization, preventing over-specialized splits and reducing variance at the cost of potential bias.
    \item \textbf{Random Forest Ensemble Size (\texttt{n\_estimators}):} Controls the number of de-correlated decision trees $B \in \{50, 100, 200\}$ aggregated via bootstrap aggregating. Leo Breiman proved that as ensemble size increases, Random Forests do not overfit; rather, generalization error converges to an upper bound:
    \begin{equation}
        \lim_{B \to \infty} PE^* = \frac{\rho(1 - s^2)}{s^2}
    \end{equation}
    where $\rho$ denotes mean tree-pair correlation and $s$ is classifier strength. However, increasing $B$ scales computational training and inference time linearly. Evaluating $B \in \{50, 100, 200\}$ measures the empirical point of diminishing accuracy returns against computational overhead.
\end{itemize}

\subsection{Theoretical Foundations of Feature Group Selection}
\label{sec:theory_feature_group}

Feature selection identifies a compact, highly informative subset of attributes $\mathcal{X}^* \subset \mathcal{X}$ that maximizes class separability while mitigating the curse of dimensionality. The Forest CoverType dataset presents an intrinsic dichotomy between two distinct attribute modalities:
\begin{enumerate}
    \item \textbf{Continuous Terrain Features ($p=10$):} Spatial and topographic variables (Elevation, Aspect, Slope, Hydrology distances, Roadway distance, Fire point distance, Hillshade indices). These continuous attributes directly capture physical geography, sunlight exposure, and moisture gradients.
    \item \textbf{Sparse Binary Indicators ($p=44$):} 4 binary Wilderness Area indicators and 40 binary Soil Type indicators. These one-hot encoded variables introduce high dimensionality and extreme sparsity ($>95\%$ zero entries).
\end{enumerate}
Comparing model performance on ``Terrain features only'' ($p=10$) against ``All features'' ($p=54$) isolates the marginal utility of geological indicators. This experiment evaluates whether continuous topographic features alone provide sufficient predictive signal, and how tree-based partitioning algorithms handle the addition of 44 sparse binary flags.

\section{Evaluation Metrics}
\label{sec:evaluation_metrics}

Performance is quantified across multiple complementary metrics:
\begin{itemize}
    \item \textbf{Overall Accuracy:} Proportion of correctly classified observations:
    \begin{equation}
        \text{Accuracy} = \frac{\sum_{k=1}^K TP_k}{N}
    \end{equation}
    \item \textbf{Macro-Averaged Precision:} Arithmetic mean of precision scores across all $K=7$ classes, giving equal weight to minority and majority classes:
    \begin{equation}
        \text{Macro Precision} = \frac{1}{K} \sum_{k=1}^K \frac{TP_k}{TP_k + FP_k}
    \end{equation}
    \item \textbf{Macro-Averaged Recall:} Arithmetic mean of class sensitivities:
    \begin{equation}
        \text{Macro Recall} = \frac{1}{K} \sum_{k=1}^K \frac{TP_k}{TP_k + FN_k}
    \end{equation}
    \item \textbf{Macro-Averaged F1-Score:} Harmonic mean of macro precision and macro recall:
    \begin{equation}
        \text{Macro F1} = \frac{1}{K} \sum_{k=1}^K \frac{2 \cdot P_k \cdot R_k}{P_k + R_k}
    \end{equation}
    \item \textbf{Sample Standard Deviation ($\sigma$):} Quantifies score dispersion across $M$ distinct random split seeds:
    \begin{equation}
        \sigma = \sqrt{\frac{1}{M-1} \sum_{m=1}^M (Acc_m - \bar{Acc})^2}
    \end{equation}
    \item \textbf{Computational Execution Time ($T_{\text{sec}}$):} Total wall-clock time in seconds recorded for model fitting and batch test inference.
\end{itemize}

\section{Summary}
\label{sec:method_summary}

The proposed methodology establishes a rigorous, controlled experimental framework for evaluating machine learning stability, capacity scaling, and feature sensitivity. By standardizing dataset subsampling, feature preprocessing, hyperparameter configurations, and evaluation metrics, the pipeline isolates pseudo-random seed effects and algorithmic parameters across distinct model families. The subsequent chapters detail the system architecture and empirical findings produced by this framework.
